\chapter{Источники и порядок чтения}

\section{Рабочий набор литературы}

\begin{center}\small
\begin{tabularx}{\textwidth}{@{}P{0.22\textwidth}P{0.2\textwidth}Y@{}}
\toprule
Источник & Когда & Материал и результат\\
\midrule
Любич~\cite{LyubichBook} & 2 ч в неделю & Выбранные разделы конформной
геометрии; определение, формулировка и один пример\\
Михайлов--Пасси~\cite{MP} & 2 ч в неделю & Коммутаторы, групповые кольца,
размерные подгруппы; доказательство на 1--2 недели\\
Хэтчер~\cite{Hatcher} & Недели 1--8 & §§1.1, 1.3; §1.2 для вычисления
фундаментальной группы проколотой плоскости\\
Милнор~\cite{MilnorThird,MilnorNotes} & Недели 9--12 и 21--28 & Fatou--Julia,
локальная динамика, заполненное множество Жюли и внешние лучи\\
Некрашевич~\cite{NekraIMG,NekraBook} & После топологического блока &
Сначала §§2--3 обзора; монография по выбранной задаче\\
\bottomrule
\end{tabularx}
\end{center}

\section{Выборочный маршрут по Любичу}
Местная версия --- \repo{refs/Lyubich_book_vol_I-II_2023.pdf}.
Пагинация сверена по оглавлению.
\begin{enumerate}
\item §7, с.~133--145: текущий ознакомительный проход и выбранные доказательства.
\item §8, с.~146--151, и §10, с.~163--172: граничное продолжение и
      необходимые леммы теории потенциала.
\item §6, с.~119--132: возвращаться для доказательств о модуле.
      §9, с.~152--162: читать при работе с puzzle после топологического блока.
\item §§11--17: определения qc-геометрии, искажение модуля, MRMT,
      голоморфные движения; 2--3 выбранных доказательства.
\item §§28, 31, 33--34, 38, 40, 42, 47: материал исследовательской ветки
      после выбора конкретного механизма.
\end{enumerate}
Вводную динамику изучать по Милнору. К соответствующим пунктам §§19--24
Любича обращаться за деталями доказательства. §18 и остальные разделы
добавлять по зависимостям задачи.

\section{Доказательства из Михайлова--Пасси}
Первые ориентиры --- §1.1, §1.2, начало главы~2 и §4.1.
Определения размерных подгрупп брать над \(\mathbb Z\):
\[
\Delta=\ker\bigl(\epsilon:\mathbb Z[G]\to\mathbb Z\bigr),\qquad
D_n(G)=\{g\in G:g-1\in\Delta^n\}.
\]
На квартал выбрать 4--6 доказательств из следующей цепочки:
\begin{enumerate}
\item Вывести коммутаторные тождества и доказать характеристичность
      \(\gamma_n(G)\). Проверить включение
\[
\gamma_n(G)/\gamma_{n+1}(G)\subseteq
\operatorname{Z}\bigl(G/\gamma_{n+1}(G)\bigr).
\]
\item Изоморфизм \(G_{\mathrm{ab}}\simeq\Delta/\Delta^2\), заданный
      \(g\mapsto(g-1)+\Delta^2\); следствие \(D_2(G)=\gamma_2(G)\).
\item Нормальность \(D_n(G)\) и включение
      \([D_m(G),D_n(G)]\subseteq D_{m+n}(G)\). Использовать тождество
\[
[x,y]-1=x^{-1}y^{-1}\bigl((x-1)(y-1)-(y-1)(x-1)\bigr),
\]
      где \([x,y]=x^{-1}y^{-1}xy\).
\item Включение \(\gamma_n(G)\subseteq D_n(G)\) по индукции.
\item Одна лемма о свободной группе или следствие коммутаторного исчисления
      из §§1.1--1.2; выписать все используемые тождества.
\item Разложение Магнуса \(x_i\mapsto1+X_i\): вычислить коммутатор
      до степени~2 в некоммутирующих переменных.
\end{enumerate}
На каждое доказательство оформить карточку и проверить конкретную группу.
В главах о гомологиях и симплициальных методах выбирать результаты по
исследовательскому паспорту.

\section{Топологический и динамический блоки}
По Хэтчеру разобрать подъём путей и гомотопий и действие фундаментальной
группы на слое накрытия. Вычислить накрытие \(z\mapsto z^2\) на \(\mathbb C^*\)
и один пример для проколотой плоскости. Затем открыть определения
итерационной монодромии в обзоре Некрашевича.

Для Милнора использовать третье издание либо местный ранний конспект
\repo{refs/9201272v1.pdf}. В раннем конспекте выбрать §§3, 6, 15--18;
в третьем издании ориентироваться на названия разделов.
Первый результат --- одно доказательство о локальной динамике и один
разобранный переход к заполненному множеству Жюли или внешним лучам.

\section{Современные статьи: по одной цепочке}
Для выбранного результата выписать: класс параметров, ограничения
комбинаторики, геометрическую оценку, используемый параметрический переход
и заключение. Разбирать одну зависимость за раз. Карта источников ---
\repo{refs/README.md} и \repo{refs/mlc_sources.json}.

\section{Обновление литературы}
Раз в месяц выделять до 30 минут на страницы авторов и новые версии только
выбранных работ. В журнал заносить изменения утверждений и их влияние
на текущую задачу. Раз в квартал обновлять карту литературы.

Интернет-ссылки проверены 6 ноября 2026 года. В тексте фиксировать
использованную версию статьи.

\begin{thebibliography}{99}\small
\bibitem{LyubichBook}
M.~Lyubich. \emph{Conformal Geometry and Dynamics of Quadratic Polynomials,
vol.~I--II}. Online manuscript; локальная версия 2023 года.
\href{https://www.math.stonybrook.edu/~mlyubich/book.pdf}{Книга};
\href{https://www.math.stonybrook.edu/~mlyubich/papers/index.html}{страница автора}.

\bibitem{MP}
R.~Mikhailov, I.~B.~S.~Passi.
\emph{Lower Central and Dimension Series of Groups}.
Lecture Notes in Mathematics 1952, Springer, 2009.
\href{https://link.springer.com/book/10.1007/978-3-540-85818-8}{Издательская страница и оглавление}.

\bibitem{DH}
A.~Douady, J.~H.~Hubbard.
\emph{Exploring the Mandelbrot Set: The Orsay Notes}.
\href{https://pi.math.cornell.edu/~hubbard/OrsayEnglish.pdf}{Авторская публикация}.

\bibitem{DudkoSurvey}
D.~Dudko. \emph{On the MLC Conjecture and the Renormalization Theory
in Complex Dynamics}. arXiv:2512.24171v1, 2025.
\href{https://arxiv.org/abs/2512.24171}{Версия препринта}.
Роль в плане: карта современной теории.

\bibitem{Satellite}
D.~Dudko, M.~Lyubich. \emph{Local connectivity of the Mandelbrot set
at some satellite parameters of bounded type}.
\href{https://arxiv.org/abs/1808.10425}{arXiv:1808.10425}.
Локальный PDF: \repo{refs/Dudko_1808.10425.pdf}.

\bibitem{Feigenbaum}
D.~Dudko, M.~Lyubich. \emph{MLC at Feigenbaum points}.
\href{https://arxiv.org/abs/2309.02107}{arXiv:2309.02107}.
Локальный PDF: \repo{refs/Dudko_2309.02107.pdf}.

\bibitem{KKL}
J.~Kahn, A.~Kapiamba, M.~Lyubich.
\emph{MLC for parabolically bounded primitive renormalization}.
\href{https://arxiv.org/abs/2606.27272}{arXiv:2606.27272v1}, 2026.
Локальный исходник: \repo{refs/arxiv-2606.27272/waves-with-Misha.tex}.

\bibitem{BDI}
L.~Bartholdi, D.~Dudko. \emph{Algorithmic aspects of branched coverings I.
Van Kampen's Theorem for bisets}.
\href{https://arxiv.org/abs/1512.08539}{arXiv:1512.08539}.
Локальный исходник: \repo{refs/arxiv-1512.08539/bc1-bisets.tex}.

\bibitem{BDIV}
L.~Bartholdi, D.~Dudko. \emph{Algorithmic aspects of branched coverings
IV/V. Expanding maps}.
\href{https://arxiv.org/abs/1610.02434}{arXiv:1610.02434}.
Локальный исходник: \repo{refs/arxiv-1610.02434/bc4-expanding.tex}.

\bibitem{DudkoPage}
D.~Dudko.
\href{https://www.math.stonybrook.edu/~ddudko/Research}{Авторский список работ}.
Источник актуальных названий и версий.

\bibitem{Hatcher}
A.~Hatcher. \emph{Algebraic Topology}. Cambridge University Press, 2002.
\href{https://pi.math.cornell.edu/~hatcher/AT/ATpage.html}{Авторская страница и свободный текст}.

\bibitem{MilnorThird}
J.~Milnor. \emph{Dynamics in One Complex Variable}. Third edition,
Annals of Mathematics Studies 160, Princeton University Press, 2006.
\href{https://www.jstor.org/stable/j.ctt7rnxn}{Издание Princeton University Press}.

\bibitem{MilnorNotes}
J.~Milnor. \emph{Dynamics in One Complex Variable: Introductory Lectures}.
Ранняя версия, revised 1991, \href{https://arxiv.org/abs/math/9201272}{arXiv:math/9201272}.
Местный файл: \repo{refs/9201272v1.pdf}.

\bibitem{NekraIMG}
V.~Nekrashevych. \emph{Iterated monodromy groups}.
В: \emph{Groups St Andrews 2009 in Bath}, vol.~1, LMS Lecture Note Series
387 (2011), pp.~41--93.
\href{https://people.tamu.edu/~nekrash/Preprints/img.pdf}{Предварительная авторская версия, 48 страниц}.

\bibitem{NekraBook}
V.~Nekrashevych. \emph{Self-similar groups}.
Mathematical Surveys and Monographs 117, American Mathematical Society, 2005.
\href{https://people.tamu.edu/~nekrash/publ.html}{Авторский список работ};
\href{https://www.math.stonybrook.edu/~mlyubich/Archive/Stock/Nekrash.pdf}{Ранняя версия монографии}.

\bibitem{KarpickeBlunt}
J.~D.~Karpicke, J.~R.~Blunt.
\emph{Retrieval practice produces more learning than elaborative studying
with concept mapping}. Science 331 (2011), no.~6018, pp.~772--775.
\href{https://learninglab.psych.purdue.edu/downloads/2011/2011_Karpicke_Blunt_Science.pdf}{Текст статьи}.
\end{thebibliography}
